\chapter{Introduction}
\label{ch:intro}

\section{Motivation}

A mid-sized hospital's network moves hundreds of thousands of small data
packets every minute: patient record lookups, internal email, video
calls, backup jobs. A handful belong to attackers --- someone probing
for unpatched servers, someone cycling password after password against
the staff portal, someone holding open thousands of slow connections to
exhaust the booking system. The hospital's security team --- a Security
Operations Centre, or SOC --- is responsible for spotting those few
malicious flows in the flood of normal traffic. A SOC analyst processes
50 to 200 alerts per shift; if the false-positive rate is too high, the
queue overflows and real incidents are missed. This operational
constraint --- alert fatigue --- is why this thesis treats recall at a
fixed false-positive budget as the operationally relevant metric, not F1
in isolation.

A modern enterprise generates millions of network flows per day, far
beyond what humans can inspect. Automated detection is therefore
essential, and increasingly that detection is performed by
machine-learning models rather than by hand-written rules. The promise
is straightforward: train a classifier on labelled flows, deploy it at
the network perimeter, and let it surface only the suspicious flows for
human triage. Public flow datasets such as
CICIDS2017~\cite{Sharafaldin2018} and UNSW-NB15~\cite{Moustafa2015}
have made it cheap to train such classifiers, and dozens of recent
papers report macro F1 scores above 0.95.

\plainbox{A signature is a rule like ``if the URL contains
\texttt{/admin.php}, raise an alert''. It works for known patterns but
breaks the moment the attacker changes one character. Machine learning
instead learns from labelled examples, the same way a child learns to
recognise dogs from pictures rather than from a checklist of features.}

\section{The Honesty Gap}

Despite these encouraging numbers, the problem is not solved. Sommer and
Paxson~\cite{Sommer2010} long ago noted that machine learning in security
operates ``outside the closed world'' of standard ML benchmarks ---
adversaries adapt, distributions shift, and the cost of a false positive
is asymmetric. More recently, Arp et al.~\cite{Arp2022} catalogued ten
common methodological pitfalls in security-ML papers, several of which
directly motivate the evaluation choices in this thesis. Practitioners
report that deployed IDS produce false-positive rates that quickly make
the system unusable, and recall on previously unseen attack variants is
far below the published numbers. The goal of this thesis is to identify
what breaks between the laboratory benchmark and the production deployment.

\section{BSc Learning Objectives}

The work fulfils the DTU Compute BSc learning objectives in five concrete
ways. It designs and implements a complete end-to-end machine-learning
pipeline. It applies statistical rigour through bootstrap confidence
intervals, McNemar's test, and explicit threshold selection via
Youden's~J. It critically appraises published benchmark numbers and
reproduces them under stronger evaluation protocols, including LODO and
cross-dataset generalisation. It assesses security-specific failure modes
through an adversarial evasion study and produces SOC-actionable output
through MITRE ATT\&CK mapping. Finally, it documents the work to a
reproducible standard with a fully scripted Makefile pipeline that any
successor can re-run end to end on a laptop in under one hour.

\section{Research Questions}

\begin{description}[leftmargin=2em, style=nextline]
    \item[RQ1 -- Temporal generalisation.] Does flow-based machine learning generalise across time? Specifically, does day-based or leave-one-day-out evaluation yield comparable performance to a random stratified split?
    \item[RQ2 -- Adversarial robustness.] Are flow-based detectors robust to low-effort evasion? Specifically, can a score-query black-box attacker move a small subset of controllable flow features within a small $L^\infty$ budget to flip predictions, and do these adversarial examples transfer between models?
    \item[RQ3 -- Operational triage.] Can model output drive triage? Specifically, can predicted attack types be mapped reliably to MITRE ATT\&CK techniques to produce alerts that a human analyst can act on?
\end{description}

\section{Main Findings at a Glance}
\label{sec:findings-dashboard}

\begin{table}[H]
    \centering
    \footnotesize
    \renewcommand{\arraystretch}{1.4}
    \begin{tabular}{|p{3.5cm}|p{10.5cm}|}
        \hline
        \textbf{Multi-class collapse} & Macro F1 falls from 0.86 (stratified XGBoost) to 0.44 (day split, all four models within 0.04). Friday's classes never appear in Mon--Wed training. \\
        \hline
        \textbf{Binary survives} & Random forest reaches F1 = 0.86 on day split and ROC-AUC = $0.926 \pm 0.053$ across LODO folds, the lowest variance of the four models. \\
        \hline
        \textbf{Adversarial brittleness} & Undefended RF flips 21.4\% of attack flows at $\varepsilon = 0.25\sigma$ (24.4\% unbounded headline); adversarial flows transfer to XGBoost at 86.4\%. Naïve adversarial training fails. \\
        \hline
        \textbf{No universal winner} & Model ranking flips between datasets: XGBoost wins CICIDS stratified multi-class, RF wins CICIDS day-split binary, LightGBM wins UNSW-NB15 multi-class. \\
        \hline
    \end{tabular}
\end{table}

\section{Scope and Delimitations}

Two public datasets: CICIDS2017 (primary) and UNSW-NB15 (cross-dataset
validation). Four classical supervised models (LogReg, RF, XGBoost,
LightGBM) plus two unsupervised baselines (Isolation Forest, LOF). Deep
neural networks are excluded because tree ensembles match or exceed deep
models on tabular flow features~\cite{Grinsztajn2022} and the additional
complexity would crowd out the adversarial and operational analyses. The
adversarial threat model is score-query black-box with
$L^\infty$-bounded perturbation on a 19-feature subset. The operating
context is offline batch scoring; no real-time deployment or streaming
evaluation is attempted.

\section{Contributions}

The originality of this thesis is not in the choice of models, the choice
of datasets, or any individual technique. It is in the combination of
three rarely-co-located components in one reproducible end-to-end
pipeline: honest temporal evaluation that exposes the gap between
published benchmarks and deployment performance, a score-query black-box
adversarial robustness analysis on a 19-feature attacker-controllable
subset, and a static MITRE ATT\&CK triage layer that turns model output
into SOC-actionable alerts. Each component exists in prior literature in
isolation; their integration into a single end-to-end study, with
explicit negative results where they appear, is the contribution of this
thesis.

Concretely, the thesis delivers: empirical evidence that headline
benchmark numbers on stratified splits over-state operational
performance by 0.3--0.4 macro F1 units; a LODO study confirming random
forest as the most temporally robust binary detector (ROC-AUC = $0.926
\pm 0.053$); a greedy evasion plus transferability plus
adversarial-training study demonstrating brittleness and the failure of
naïve defences; a static MITRE ATT\&CK mapping routed through Boyd's
OODA loop for SOC integration; and a scripted, reproducible pipeline
that produces every figure, table, and statistical test from raw CSVs
in 45--60 minutes.

\section{Thesis Structure}

Five parts. Part~I (Foundations): introduction and theoretical
background (Chapter~\ref{ch:background}). Part~II (Method): datasets
(Chapter~\ref{ch:datasets}) and experimental methodology
(Chapter~\ref{ch:method}). Part~III (Results): multi-class
(Chapter~\ref{ch:multiclass}), binary (Chapter~\ref{ch:binary}), LODO
(Chapter~\ref{ch:lodo}), cross-dataset (Chapter~\ref{ch:crossdataset}),
adversarial (Chapter~\ref{ch:adversarial}), and best-model selection
(Chapter~\ref{ch:bestmodel}). Part~IV (Practice): operational triage
including the OODA loop and MITRE mapping (Chapter~\ref{ch:mitre}),
and deployment considerations (Chapter~\ref{ch:deployment}). Part~V
(Closing): discussion and limitations (Chapter~\ref{ch:discussion}),
and conclusion (Chapter~\ref{ch:conclusion}).
